\def\GraphicsFolder{appendices/mean-field-hubbard/pictures}
\chapter{Mean-Field Theory in Hubbard lattices}\label{appendix:mean-field-hubbard}

In this Appendix the MFT solutions to the conventional Hubbard hamiltonian of Eq.~\eqref{eq:hubbard-model},
\begin{equation}\label{appeq:hubbard-model}
	\hat H = 
	-t \sum_{\langle ij \rangle} \sum_\sigma \left[
		\hat c_{i\sigma}^\dagger \hat c_{j\sigma} + \hc
	\right]
	+ U \sum_i \hat n_{i\uparrow} \hat n_{i\downarrow}
	\qquad
	t, U  > 0
\end{equation}
are described. The discussion is limited to the two-dimensional square lattice. The results of ths appendix are relevant and used throughout the text as the groundwork to be extended by the means of MFT applied to the non-local part of the EHM.

\begin{figure}
	\centering
%	\input{\GraphicsFolder/ferromagnetic-3d-band}
	\caption{Depiction of the Hubbard square lattice hopping band $\epsilon_{\mathbf{k}} = -2t[\cos(k_x) + \cos(k_y)]$. The red lines mark the zero-energy intersection.}
	\label{appfig:ferromagnetic-3d-band}
\end{figure}

\section{Antiferromagnetic phase}\label{appsec:antiferromagnetic-phase}

We work in a commensurate AF phase ordering, which is, we explicitly break symmetry by realizing a SDW phase with spin density unevenly distributed on sites with periodicity of two sites, each direction. Recalling the discussion of Sec.~\ref{sec:ehm-symmetries}, AF long-range ordering we deal with here lies within the $\mathrm{SU}(2) \to \mathrm{U}^z(1)$ SSB, represented pictorially in Fig.~\ref{fig:su2-u1-domains}. Recalling Tab.~\ref{tab:U-wick-terms-phases}, for such a physical phase the only Wick contraction channel available is the Hartree one. Let us proceed in such sense.

\subsection{MFT treatment of the Hartree channel}\label{appsubsec:af-hartree-channel}

Let us change notation, indicating each site as
\[
	i \to \mathbf{r} = (x,y)
	\qquad
	x,y \in \mathbb{N}
\]
In terms of the physical magnetization $m$, the MFT Ansatz is given by
\begin{equation}\label{appeq:af-mean-field-ansatz}
	\langle \hat n_{\mathbf{r}\sigma} \rangle = n - (-1)^{x+y+\delta_{\sigma=\uparrow}} m
\end{equation}
Decoupling the HM hamiltonian and substituting,
\begin{align}
	\hat H &\simeq -t \sum_{\langle \mathbf{r} \mathbf{r}' \rangle} \left[
		\hat c_{i\sigma}^\dagger \hat c_{j\sigma} + \hc
	\right] + U \sum_\mathbf{r} \left[
		\hat n_{\mathbf{r}\uparrow} \ev{\hat n_{\mathbf{r}\downarrow}} + \ev{\hat n_{\mathbf{r}\uparrow}} \hat n_{\mathbf{r}\downarrow} - \ev{\hat n_{\mathbf{r}\uparrow}} \ev{\hat n_{\mathbf{r}\downarrow}}
	\right] \nonumber \\
	&= -t \sum_{\langle \mathbf{r}\mathbf{r}' \rangle} \sum_\sigma \left[
		\hat c_{i\sigma}^\dagger \hat c_{j\sigma} + \hc
	\right]
	+ nU \sum_\mathbf{r} \left[
		\hat n_{\mathbf{r}\uparrow} + \hat n_{\mathbf{r}\downarrow}
	\right] - mU \sum_\mathbf{r} (-1)^{x+y} \left[
		\hat n_{\mathbf{r}\uparrow} - \hat n_{\mathbf{r}\downarrow}
	\right] - E_{\mathrm{H}/U}^{(\mathrm{AF})} \label{appeq:af-hm-mft-h}
\end{align}
with $E_{\mathrm{H}/U}^{(\mathrm{AF})}$ the HM Hartree contraction energy,
\begin{align}
	E_{\mathrm{H}/U}^{(\mathrm{AF})} &\equiv U \sum_\mathbf{r} \ev{\hat n_{\mathbf{r}\uparrow}} \ev{\hat n_{\mathbf{r}\downarrow}} \nonumber \\
	&= U \sum_\mathbf{r} \left(
		n + (-1)^{x+y} m
	\right) \left(
		n - (-1)^{x+y} m
	\right) \nonumber \\
	&= L_x L_y \times U \left(
		n^2 - m^2
	\right) \label{appeq:af-hm-hartree-U-Ec}
\end{align}
By a Fourier transform, the phase factor of Eq.~\eqref{appeq:af-hm-mft-h} can be absorbed in the destruction operator inside of $\hat n_{\mathbf{r}\sigma}$:
\begin{align}
	\sum_\mathbf{r} (-1)^{x+y} \hat n_{\mathbf{r}\sigma} &= \sum_\mathbf{r} (-1)^{x+y} \hat c_{\mathbf{r}\sigma}^\dagger \hat c_{\mathbf{r}\sigma} \nonumber \\
	&\stackrel{\mathcal{F}}{=} \sum_\mathbf{r} e^{i \bm{\pi} \cdot \mathbf{r}}
	\frac{1}{\sqrt{L_x L_y}} \sum_{\mathbf{k} \in \mathrm{BZ}} e^{i \mathbf{k} \cdot \mathbf{r}} \hat c_{\mathbf{k}\sigma}^\dagger \frac{1}{\sqrt{L_x L_y}} \sum_{\mathbf{k}' \in \mathrm{BZ}} e^{-i \mathbf{k}' \cdot \mathbf{r}} \hat c_{\mathbf{k}'\sigma} \nonumber \\
	&=\sum_{\mathbf{k} \in \mathrm{BZ}} \sum_{\mathbf{k}' \in \mathrm{BZ}} \hat c_{\mathbf{k}\sigma}^\dagger \hat c_{\mathbf{k}'\sigma} \frac{1}{L_x L_y} \sum_\mathbf{r} e^{-i [\mathbf{k}' - (\mathbf{k} + \bm{\pi}) ] \cdot \mathbf{r}} \nonumber \\
	&= \sum_{\mathbf{k} \in \mathrm{BZ}} \hat c_{\mathbf{k}\sigma}^\dagger \hat c_{\mathbf{k}+\bm{\pi}\sigma}\label{appeq:af-hm-staggered-n-transform}
\end{align}
where $\bm{\pi} = (\pi,\pi)$. It follows:
\[
mU \sum_\mathbf{r} (-1)^{x+y} \left[
		\hat n_{\mathbf{r}\uparrow} - \hat n_{\mathbf{r}\downarrow}
	\right] = mU \sum_{\mathbf{k} \in \mathrm{BZ}} \left[
		\hat c_{\mathbf{k}\uparrow}^\dagger \hat c_{\mathbf{k}+\bm{\pi}\uparrow} - \hat c_{\mathbf{k}\downarrow}^\dagger \hat c_{\mathbf{k}+\bm{\pi}\downarrow}
	\right]
\]
Note that the Hartree channel renormalizes chemical potential,
\[
	nU \sum_\mathbf{r} \left[
	\hat n_{\mathbf{r}\uparrow} + \hat n_{\mathbf{r}\downarrow}
	\right] = nU \times \hat N
	\quad\implies\quad
	\mu \to \tilde{\mu} \equiv \mu - nU
\]
an effect that is present also in the EHM.

\begin{figure}
	\centering
%	\subfloat[Alternative band.]{
%		\input{\GraphicsFolder/alternative-ferromagnetic-3d-band}
%		\label{appsubfig:alternative-ferromagnetic-3d-band}
%	}
%	\subfloat[Contour plot.]{
%		\input{\GraphicsFolder/ferromagnetic-contour}
%		\label{appsubfig:ferromagnetic-contour}
%	}
	\caption{Alternative depiction of the Hubbard square lattice hopping band previously reported in Fig.~\ref{appfig:ferromagnetic-3d-band}. The Magnetic Brillouin Zone ($\mathrm{MBZ}$) is delimited by the zero-energy contour and is indicated in figure. As it is evident, energy sign flips by taking a $(\pi,\pi)$ translation in $\mathbf{k}$ space.}
	\label{appfig:alternative-ferromagnetic-band}
\end{figure}

\subsection{Nambu representation of the AF phase and diagonalization}\label{appsubsec:af-nambu}

Consider the band of Fig.~\ref{appfig:ferromagnetic-3d-band} . As does \citeauthor{fabrizio2022course} \cite{fabrizio2022course}, the area delimited externally by the solid lines at zero energy is denominated ``Magnetic Brillouin Zone'' ($\mathrm{MBZ}$). The periodicity of $\mathbf{k}$ space guarantees that the full $\mathrm{BZ}$ can be taken as well to be the one of Fig.~\ref{appsubfig:alternative-ferromagnetic-3d-band}. Then:
\begin{align}
	\sum_{\mathbf{k} \in \mathrm{BZ}}
	\hat c_{\mathbf{k}\sigma}^\dagger \hat c_{\mathbf{k}+\bm{\pi}\sigma} &= \sum_{\mathbf{k} \in \mathrm{MBZ}}
	\left[
		\hat c_{\mathbf{k}\sigma}^\dagger \hat c_{\mathbf{k}+\bm{\pi}\sigma} + \hat c_{\mathbf{k}+\bm{\pi}\sigma}^\dagger \hat c_{\mathbf{k}+2\bm{\pi}\sigma}
	\right] \nonumber \\
	&= \sum_{\mathbf{k} \in \mathrm{MBZ}}
	\left[
		\hat c_{\mathbf{k}\sigma}^\dagger \hat c_{\mathbf{k}+\bm{\pi}\sigma} + \hat c_{\mathbf{k}+\bm{\pi}\sigma}^\dagger \hat c_{\mathbf{k}\sigma}
	\right] \label{appeq:mixed-cc-term-intermediate-passage}
\end{align}
Periodicity by shifts $2\bm{\pi}$ has been used. Now, define the Nambu spinors:
\[
	\hat \Psi_{\mathbf{k}\sigma} \equiv \begin{bmatrix}
		\hat c_{\mathbf{k}\sigma} \\
		\hat c_{\mathbf{k}+\bm{\pi}\sigma} 
	\end{bmatrix}
\]
As in Eq.~\mref, it is useful to define the density and AF fluctuation operators:
\begin{align}
	\hat n_{\mathbf{k}\sigma} &\equiv \hat c_{\mathbf{k}\sigma}^\dagger \hat c_{\mathbf{k}\sigma} \label{appeq:af-hm-n-def} \\
	\hat \psi_{\mathbf{k}\sigma} &\equiv \hat c_{\mathbf{k}\sigma}^\dagger \hat c_{\mathbf{k}+\bm{\pi}\sigma} \label{appeq:af-hm-psi-def}
\end{align}
whose EVs are relevant order parameters for the normal and AF phases. Let also the gap function:
\[
	\Delta_{\mathbf{k}\sigma} = (-1)^{\delta_{\sigma=\downarrow}} mU
\]

At fixed filling, the $nU$ term is a pure chemical potential shift, thus will be neglected. The kinetic term transforms as
\[
\begin{aligned}
	-t \sum_{\langle \mathbf{r}\mathbf{r}' \rangle} \sum_\sigma \left[
		\hat c_{i\sigma}^\dagger \hat c_{j\sigma} + \hc
	\right] &\stackrel{\mathcal{F}}{=} -t \sum_{\langle \mathbf{r}\mathbf{r}' \rangle} \sum_\sigma \frac{1}{\sqrt{L_xL_y}} \sum_{\mathbf{k}\in\mathrm{BZ}} e^{i\mathbf{k}\cdot\mathbf{r}} \hat c_{\mathbf{k}\sigma}^\dagger \frac{1}{\sqrt{L_xL_y}} \sum_{\mathbf{k}'\in\mathrm{BZ}} e^{-i\mathbf{k}'\cdot\mathbf{r}'} \hat c_{\mathbf{k}'\sigma} +\hc \\
	&= -2t\sum_{\mathbf{k},\mathbf{k}'} \sum_\sigma \hat c_{\mathbf{k}\sigma}^\dagger \hat c_{\mathbf{k}'\sigma} \frac{1}{L_x L_y} \sum_{\mathbf{r}\in\mathcal{S}/2} e^{i(\mathbf{k}-\mathbf{k}')\cdot\mathbf{r}} \left[
		e^{i k'_x}+e^{-i k'_x} + e^{i k'_y}+e^{-i k'_y}
	\right]  \\
	&= \sum_{\mathbf{k} \in \mathrm{BZ}} \sum_\sigma \epsilon_{\mathbf{k}\sigma} \hat c_{\mathbf{k}\sigma}^\dagger \hat c_{\mathbf{k}\sigma} \\ 
	&= \sum_{\mathbf{k} \in \mathrm{MBZ}} \sum_\sigma \left[
		\epsilon_{\mathbf{k}\sigma} \hat c_{\mathbf{k}\sigma}^\dagger \hat c_{\mathbf{k}\sigma}
		+ \epsilon_{\mathbf{k}+\bm{\pi}\sigma} \hat c_{\mathbf{k}+\bm{\pi}\sigma}^\dagger \hat c_{\mathbf{k}+\bm{\pi}\sigma}
	\right] \\
	&= \sum_{\mathbf{k} \in \mathrm{MBZ}} \sum_\sigma \epsilon_{\mathbf{k}\sigma}
	\left[
		\hat c_{\mathbf{k}\sigma}^\dagger \hat c_{\mathbf{k}\sigma}
		- \hat c_{\mathbf{k}+\bm{\pi}\sigma}^\dagger \hat c_{\mathbf{k}+\bm{\pi}\sigma}
	\right] \\ 
\end{aligned}
\]
where $\epsilon_{\mathbf{k}\sigma} \equiv -2t ( \cos k_x + \cos k_y )$. In the fourth passage, the sum over the full $\mathrm{BZ}$ was written considering that the entirety of the zone is given by all the points in the $\mathrm{MBZ}$ plus their conjugates obtained by a $\bm{\pi}$ shift in the flipped band. As depicted in Fig.~\ref{appsubfig:alternative-ferromagnetic-3d-band}, kinetic energy is anti-periodic in $\mathbf{k}$ space by a vector $\bm{\pi}$. This anti-periodicity accounts for the minus sign arising in the last passage. The hamiltonian is then given by:
\begin{equation}\label{af-hm-bogoliubov-h}
	\hat H \simeq \sum_{\mathbf{k} \in \mathrm{MBZ}} \sum_\sigma \hat \Psi_{\mathbf{k}\sigma}^\dagger h_{\mathbf{k}\sigma} \hat \Psi_{\mathbf{k}\sigma} - E_{\mathrm{H}/U}^{(\mathrm{AF})}
	\qq{being}
	h_{\mathbf{k}\sigma} \equiv \begin{bmatrix}
		\epsilon_{\mathbf{k}\sigma} & -\Delta_{\mathbf{k}\sigma} \\
		-\Delta_{\mathbf{k}\sigma}^* & - \epsilon_{\mathbf{k}\sigma}
	\end{bmatrix}
\end{equation}
Notice: the Nambu hamiltonian is a $2\times2$ matrix over the $\mathrm{MBZ}$ -- which is half the full $\mathrm{BZ}$, coherently with a solution which essentially bipartites the lattice giving back a double sized unit cell.

\begin{figure}
	\centering
	% \input{\GraphicsFolder/pseudo-magnetic-field}
	\caption{Pseudo-magnetic field originating from mean-field treatment of the square Hubbard hamiltonian. Here, only the $\sigma=\uparrow$ situation is plotted.}
	\label{appfig:pseudo-magnetic-field}
\end{figure}

The system ground-state is obtained by the means of a Bogoliubov rotation. The hamiltonian maps onto the simple one of a spin in a magnetic field,
\[
	h_{\mathbf{k}\sigma} = \epsilon_{\mathbf{k}\sigma} \tau^z - \Delta_{\mathbf{k}\sigma} \tau^x
\]
being $\tau^\alpha$ the Pauli matrices. Then, defining
\[
	\hat s_{\mathbf{k}\sigma}^\alpha \equiv \hat \Psi_{\mathbf{k}\sigma}^\dagger \tau^\alpha \hat \Psi_{\mathbf{k}\sigma}
	\qq{and}
	\mathbf{b}_{\mathbf{k}\sigma} \equiv \begin{bmatrix}
		-\Delta_{\mathbf{k}\sigma} \\ 0 \\ \epsilon_{\mathbf{k}\sigma} 
	\end{bmatrix}
\]
one gets:
\begin{equation}\label{af-hm-spin-h}
	\hat H \simeq \sum_{\mathbf{k} \in \mathrm{MBZ}} \sum_\sigma \mathbf{b}_{\mathbf{k}\sigma} \cdot \hat{\mathbf{s}}_{\mathbf{k}\sigma} - E_{\mathrm{H}/U}^{(\mathrm{AF})}
	\qq{where}
	\hat{\mathbf{s}}_{\mathbf{k}\sigma} = \begin{bmatrix}
		\hat s_{\mathbf{k}\sigma}^x \\
		\hat s_{\mathbf{k}\sigma}^y \\
		\hat s_{\mathbf{k}\sigma}^z
	\end{bmatrix}
\end{equation}
The hamiltonian represents a system of spins subject to local magnetic fields, each tilted by an angle $\tan(2\theta_{\mathbf{k}\sigma}) = \Delta_{\mathbf{k}\sigma}/\epsilon_{\mathbf{k}\sigma}$, as sketched in Fig.~\ref{appfig:pseudo-magnetic-field}. Evidently the following relations, also reported in Tab.~\ref{tab:af-pseudospin-map}, hold:
\begin{align}
	\hat s_{\mathbf{k}\sigma}^x &= \hat \psi_{\mathbf{k}\sigma} + \hat \psi_{\mathbf{k}\sigma}^\dagger \label{appeq:af-hm-x-psi} \\
	\hat s_{\mathbf{k}\sigma}^y &= -i \left(
		\hat \psi_{\mathbf{k}\sigma} - \hat \psi_{\mathbf{k}\sigma}^\dagger 
	\right) \label{appeq:af-hm-y-psi} \\
	\hat s_{\mathbf{k}\sigma}^z &= \hat n_{\mathbf{k}\sigma} - \hat n_{\mathbf{k}+\bm{\pi}\sigma} \label{appeq:af-hm-z-psi} 
\end{align}
Let also:
\begin{align}
	\hat \Id_{\mathbf{k}\sigma} &\equiv \hat \Psi_{\mathbf{k}\sigma}^\dagger \Id_{2\times2} \hat \Psi_{\mathbf{k}\sigma} \nonumber \\
	&= \hat n_{\mathbf{k}\sigma} + \hat n_{\mathbf{k}+\bm{\pi}\sigma} \label{appeq:af-hm-id-def}
\end{align}
which accounts for a chemical potential shift.

Diagonalization of each $h_{\mathbf{k}\sigma}$ is obtained trivially by a simple rotation around the $y$ axis:
\[
	d_{\mathbf{k}\sigma} = W_{\mathbf{k}\sigma} h_{\mathbf{k}\sigma} W_{\mathbf{k}\sigma}^\dagger
	\qq{with}
	d_{\mathbf{k}\sigma} \equiv \begin{bmatrix}
		E_{\mathbf{k}\sigma} & \\
		& - E_{\mathbf{k}\sigma}
	\end{bmatrix}
\]
and where the diagonalization unitary matrix is simply a rotation by an angle $2 \theta_{\mathbf{k}\sigma}$ around the $y$ axis, counter-clockwise
\begin{align}
	W_{\mathbf{k}\sigma} &\equiv e^{-i \theta_{\mathbf{k}\sigma} \tau^y} \nonumber\\
	&= \cos \theta_{\mathbf{k}\sigma} -i \sin \theta_{\mathbf{k}\sigma} \tau^y \nonumber\\
	&= \begin{bmatrix}
		\cos \theta_{\mathbf{k}\sigma} & - \sin \theta_{\mathbf{k}\sigma} \\
		\sin \theta_{\mathbf{k}\sigma} & \cos \theta_{\mathbf{k}\sigma}
	\end{bmatrix} \label{appeq:af-hm-diagonalizer}
\end{align}
Eigenvalues are:
\begin{equation}\label{appeq:af-bands}
	E_{\mathbf{k}\sigma} \equiv \sqrt{
		\epsilon_{\mathbf{k}\sigma}^2 + \abs{\Delta_{\mathbf{k}\sigma}}^2
	}
\end{equation}

At any finite temperature, the ground-state of such a system is well-known: a thermal superposition of ``down'' spins with ``up'' spins. In the pseudospin picture, at each ``site'' $(\mathbf{k}\sigma)$, the lowest-energy state is the one anti-aligned with the pseudo-magnetic field. It is useful now to define the thermal factors,
\begin{equation}\label{appeq:af-thermal-factors-def}
	\mathrm{Th}_{\mathbf{k}\sigma}^{(\pm)}[\beta,\mu] \equiv 
	f\left(
		E_{\mathbf{k}\sigma};\beta,\mu
	\right) \pm f\left(
		-E_{\mathbf{k}\sigma};\beta,\mu
	\right)
\end{equation}
that will appear quite often. Then it's immediate to derive the various expectation values.
\begin{align}
	\ev{\hat s_{\mathbf{k}\sigma}^x} &= -\frac{\Re\lbrace\Delta_{\mathbf{k}\sigma}\rbrace}{2E_{\mathbf{k}\sigma}} \mathrm{Th}_{\mathbf{k}\sigma}^{(-)}[\beta,\mu] \label{appeq:af-hm-sx}\\
	\ev{\hat s_{\mathbf{k}\sigma}^y} &= -\frac{\Im\lbrace\Delta_{\mathbf{k}\sigma}\rbrace}{2E_{\mathbf{k}\sigma}}  \mathrm{Th}_{\mathbf{k}\sigma}^{(-)}[\beta,\mu] \label{appeq:af-hm-sy}\\
	\ev{\hat s_{\mathbf{k}\sigma}^z} &=  \frac{\epsilon_{\mathbf{k}\sigma}}{2E_{\mathbf{k}\sigma}}  \mathrm{Th}_{\mathbf{k}\sigma}^{(-)}[\beta,\mu] \label{appeq:af-hm-sz}
\end{align}

%Note that rotation is taken to be of an angle $2\theta_{\mathbf{k}\sigma} - \pi$ in order to anti-align the pseudo-spin with the magnetic field of Fig.~\ref{appfig:pseudo-magnetic-field} and thus order the eigenvalues as in $d_{\mathbf{k}\sigma}$, with the smaller one in the top-left corner of the matrix and the larger one in the bottom-right corner. The explicit form of the transformation matrix is thus
%\begin{equation}\label{appeq:W-explicit-form}
%	W_{\mathbf{k}\sigma} = \begin{bmatrix}
%		\cos \theta_{\mathbf{k}\sigma} & - \sin \theta_{\mathbf{k}\sigma} \\
%		\sin \theta_{\mathbf{k}\sigma} & \cos \theta_{\mathbf{k}\sigma}
%	\end{bmatrix} \begin{bmatrix}
%		& -1 \\
%		1 &
%	\end{bmatrix} = \begin{bmatrix}
%		- \sin \theta_{\mathbf{k}\sigma} & - \cos \theta_{\mathbf{k}\sigma} \\ 
%		\cos \theta_{\mathbf{k}\sigma} & - \sin \theta_{\mathbf{k}\sigma}
%	\end{bmatrix}
%\end{equation}


Note that considering the gran canonical hamiltonian $\hat K \equiv \hat H - \mu\hat N$ and using Eq.~\eqref{appeq:af-hm-id-def} the diagonalizing matrix remains the same, being
\begin{equation}\label{appeq:af-full-h-diagonalization}
	-\mu\hat N = \sum_{\mathbf{k}\in\mathrm{MBZ}} \begin{bmatrix}
		-\mu & \\ & -\mu
	\end{bmatrix}
	\quad\implies\quad
	W_{\mathbf{k}\sigma} \left[ 
	h_{\mathbf{k}\sigma}-\mu \Id_{2\times2}
	\right] W_{\mathbf{k}\sigma}^\dagger = d_{\mathbf{k}\sigma}-\mu \Id_{2\times2}
\end{equation}
thus for non-null chemical potential (i.e. finite doping) the eigenvalues are just shifted, $\pm E_\mathbf{k}-\mu$. In other words, the new bands retain the same chemical potential of the bare bands. Notice also that the presence of an absolute value makes the eigenvalues independent of the $\sigma$ index. Eigenvector fermion operators are obtained simply as:
\begin{equation}\label{appeq:af-hm-Gamma-WPsi}
	\hat\Gamma_{\mathbf{k}\sigma} = W_{\mathbf{k}\sigma} \hat \Psi_{\mathbf{k}\sigma}
\end{equation}
The $\hat\Gamma$ spinor operators are effectively free fermionic spinor fields and as such must be treated. Explicitly,
\begin{align}
	\hat\Gamma_{\mathbf{k}\sigma} &= \begin{bmatrix}
		\cos \theta_{\mathbf{k}\sigma} & - \sin \theta_{\mathbf{k}\sigma} \\
		\sin \theta_{\mathbf{k}\sigma} & \cos \theta_{\mathbf{k}\sigma}
	\end{bmatrix} \begin{bmatrix}
		\hat c_{\mathbf{k}\sigma} \\
		\hat c_{\mathbf{k}+\bm{\pi}\sigma} 
	\end{bmatrix} \nonumber \\
	&= \begin{bmatrix}
		\cos \theta_{\mathbf{k}\sigma} \hat c_{\mathbf{k}\sigma} - \sin \theta_{\mathbf{k}\sigma} \hat c_{\mathbf{k}+\bm{\pi}\sigma}  \\
		\sin \theta_{\mathbf{k}\sigma} \hat c_{\mathbf{k}\sigma} + \cos \theta_{\mathbf{k}\sigma} \hat c_{\mathbf{k}+\bm{\pi}\sigma} 
	\end{bmatrix} \equiv
	\begin{bmatrix}
		\hat \gamma_{\mathbf{k}\sigma}^{(+)} \\
		\hat \gamma_{\mathbf{k}\sigma}^{(-)} 
	\end{bmatrix} \label{appeq:af-hm-gamma-ops-transform}
\end{align}

\begin{figure}
	\centering
	% \input{\GraphicsFolder/antiferromagnetic-3d-bands}
	\caption{Depiction of the AF bands $\pm E_\mathbf{k}$. \todo}
	\label{appfig:antiferromagnetic-3d-bands}
\end{figure}

The consequence of Eq.~\eqref{appeq:af-full-h-diagonalization} is interesting by itself: within this HF solution, the gap opens always at $E=0$. In other words, the bands sketched in Fig.~\ref{appfig:antiferromagnetic-3d-bands} retain the same geometric features regardless of the doping, what changes is the bands width and the gap between them. The very key point here is that, being the two bands symmetric (thus containing each half of the electronic states) and defined over the MBZ, any finite doping $\delta>0$ will produce a metallic behaviour, imposing that the chemical potential crosses the upper band. It is widely observed and confirmed by more refined numerical procedures, however, that AF establishes an insulating phase -- here only reproduced at half-filling, where the chemical potential lies within the gap. As will be detailed in \ref{appsubsec:instability-af-commensurate}, the commensurate nature of AF SSB at wavevector $\bm{\pi}$ of this self-consistent solution is highly unstable and atypical at finite doping \cite{singh1990collective}.

\subsection{Explicit zero-temperature AF ground state at half-filling}\label{appsubsec:af-zero-T-half-filling-gs}

The AF ground state can be expressed easily at $n=1/2$ and $\beta \to +\infty$. Let $\ket{\Uparrow_{\mathbf{k}\sigma}}$ and $\ket{\Downarrow_{\mathbf{k}\sigma}}$ be respectively the ``up'' and ``down'' $\mathbf{k}$-th pseudospin states. Let $\ket{\overline{\Downarrow}_{\mathbf{k}\sigma}}$ be the ``down'' state in the tilted frame. The ground state is given by
\begin{equation}\label{appeq:af-hm-gs-1}
	\ket{\Psi} \equiv \bigotimes_\sigma \bigotimes_{\mathbf{k}\in\mathrm{MBZ}} \ket{\overline{\Downarrow}_{\mathbf{k}\sigma}}
\end{equation}
We can link the ``down'' states via an inverse rotation,
\begin{align}
	\ket{\overline{\Downarrow}_{\mathbf{k}\sigma}} &= e^{-i\theta_{\mathbf{k}\sigma} \hat s_{\mathbf{k}\sigma}^y} \ket{\Downarrow_{\mathbf{k}\sigma}} \nonumber \\
	&= \cos \theta_{\mathbf{k}\sigma} \ket{\Downarrow_{\mathbf{k}\sigma}} - i \sin\theta_{\mathbf{k}\sigma} \hat s_{\mathbf{k}\sigma}^y \ket{\Downarrow_{\mathbf{k}\sigma}} \nonumber \\
	&= \cos \theta_{\mathbf{k}\sigma} \ket{\Downarrow_{\mathbf{k}\sigma}} - \sin\theta_{\mathbf{k}\sigma} (\hat s_{\mathbf{k}\sigma}^+ - \hat s_{\mathbf{k}\sigma}^-) \ket{\Downarrow_{\mathbf{k}\sigma}} \nonumber \\
	&= \cos \theta_{\mathbf{k}\sigma} \ket{\Downarrow_{\mathbf{k}\sigma}} - \sin\theta_{\mathbf{k}\sigma} \ket{\Uparrow_{\mathbf{k}\sigma}} \label{appeq:af-hm-tilted-down}
\end{align}
having we used
\[
	\hat s_{\mathbf{k}\sigma}^\pm \equiv \frac{\hat s_{\mathbf{k}\sigma}^x \pm i\hat s_{\mathbf{k}\sigma}^y}{2}
\]

Use now Eqns.~\eqref{appeq:af-hm-id-def} and \eqref{appeq:af-hm-z-psi},
\[
	\frac{\hat \Id_{\mathbf{k}\sigma} + \hat s_{\mathbf{k}\sigma}^z}{2} = \hat n_{\mathbf{k}\sigma}
	\qquad
	\frac{\hat \Id_{\mathbf{k}\sigma} - \hat s_{\mathbf{k}\sigma}^z}{2} = \hat n_{\mathbf{k}+\bm{\pi}\sigma}
\]
Evaluating the above operators on the ``up'' and ``down'' states, we get immediately
\[
\begin{aligned}
	\frac{1}{2} \mel{\Uparrow_{\mathbf{k}\sigma}}{\hat \Id_{\mathbf{k}\sigma} + \hat s_{\mathbf{k}\sigma}^z}{\Uparrow_{\mathbf{k}\sigma}} &= 1 \\
	\frac{1}{2} \mel{\Downarrow_{\mathbf{k}\sigma}}{\hat \Id_{\mathbf{k}\sigma} + \hat s_{\mathbf{k}\sigma}^z}{\Downarrow_{\mathbf{k}\sigma}} &= 0
\end{aligned}
\qquad
\begin{aligned}
	\frac{1}{2} \mel{\Uparrow_{\mathbf{k}\sigma}}{\hat \Id_{\mathbf{k}\sigma} - \hat s_{\mathbf{k}\sigma}^z}{\Uparrow_{\mathbf{k}\sigma}} &= 0 \\
	\frac{1}{2} \mel{\Downarrow_{\mathbf{k}\sigma}}{\hat \Id_{\mathbf{k}\sigma} - \hat s_{\mathbf{k}\sigma}^z}{\Downarrow_{\mathbf{k}\sigma}} &= 1
\end{aligned}
\]
Thus, up to a phase
\[
	\ket{\Uparrow_{\mathbf{k}\sigma}} \sim \hat c_{\mathbf{k}\sigma}^\dagger \ket{\Omega_{\mathbf{k}\sigma}}
	\qquad
	\ket{\Downarrow_{\mathbf{k}\sigma}} \sim \hat c_{\mathbf{k}+\bm{\pi}\sigma}^\dagger \ket{\Omega_{\mathbf{k}\sigma}}
\]
being $\ket{\Omega_{\mathbf{k}\sigma}}$ the $(\mathbf{k},\sigma)$ Hilbert subspace vacuum. 

It follows that composing Eqns.~\eqref{appeq:af-hm-gs-1} and \eqref{appeq:af-hm-tilted-down} the ground state $\ket{\Psi}$ can be written as:
\begin{equation}\label{appeq:af-hm-gs-2}
	\ket{\Psi} = \bigotimes_\sigma \bigotimes_{\mathbf{k}\in\mathrm{MBZ}} \left[
		\sin\theta_{\mathbf{k}\sigma} \hat c_{\mathbf{k}\sigma}^\dagger + e^{2i\zeta_{\mathbf{k}\sigma}} \cos \theta_{\mathbf{k}\sigma} \hat c_{\mathbf{k}+\bm{\pi}\sigma}^\dagger
	\right] \ket{\Omega_{\mathbf{k}\sigma}}
\end{equation}
having we absorbed in the relative phase $e^{2i\zeta_{\mathbf{k}\sigma}}$ the two states phases differences as well as the original sign difference of Eq.~\eqref{appeq:af-hm-tilted-down}. Now, inserting the explicit form of the ground state from equation above:
\[
\begin{aligned}
	\mel{\overline{\Downarrow}_{\mathbf{k}\sigma}}{\hat \psi_{\mathbf{k}\sigma} + \hat \psi_{\mathbf{k}\sigma}^\dagger}{\overline{\Downarrow}_{\mathbf{k}\sigma}} &= \left(e^{2i\zeta_{\mathbf{k}\sigma}} + e^{-2i\zeta_{\mathbf{k}\sigma}}\right) \sin\theta_{\mathbf{k}\sigma} \cos\theta_{\mathbf{k}\sigma} \\
	&= \cos 2\zeta_{\mathbf{k}\sigma} \sin 2\theta_{\mathbf{k}\sigma}
\end{aligned}
\]
but from Eq.~\eqref{appeq:af-hm-x-psi}
\[
\begin{aligned}
	\mel{\overline{\Downarrow}_{\mathbf{k}\sigma}}{\hat \psi_{\mathbf{k}\sigma} + \hat \psi_{\mathbf{k}\sigma}^\dagger}{\overline{\Downarrow}_{\mathbf{k}\sigma}} &= 	\mel{\overline{\Downarrow}_{\mathbf{k}\sigma}}{\hat s_{\mathbf{k}\sigma}^x}{\overline{\Downarrow}_{\mathbf{k}\sigma}} \\
	&= \sin 2\theta_{\mathbf{k}\sigma}
\end{aligned}
\]
since the pseudofield has no $y$ component. This implies necessarily
\[
	\cos 2\zeta_{\mathbf{k}\sigma}
	\quad\implies\quad
	\zeta_{\mathbf{k}\sigma} = k\pi
	\qq{with}
	k\in\mathbb{Z}
\]
and lets us conclude from Eq.~\eqref{appeq:af-hm-gs-2}
\begin{equation}\label{appeq:af-hm-gs-3}
	\ket{\Psi} = \bigotimes_\sigma \bigotimes_{\mathbf{k}\in\mathrm{MBZ}} \left[
	\sin \theta_{\mathbf{k}\sigma} \hat c_{\mathbf{k}\sigma}^\dagger + \cos \theta_{\mathbf{k}\sigma} \hat c_{\mathbf{k}+\bm{\pi}\sigma}^\dagger
	\right] \ket{\Omega_{\mathbf{k}\sigma}}
\end{equation}
For low-temperature physics, we can expect the ground-state density matrix to be highly dominated by this BCS-like zero-temperature ground-state, which highlights the essential features of the system. Finally, using Eq.~\eqref{appeq:af-hm-gamma-ops-transform}, we check that correctly the ground-state describes a perfect degenerate $\gamma^{(-)}$ quasiparticles Fermi gas,
\begin{equation}\label{appeq:af-hm-gs-4}
	\ket{\Psi} = \bigotimes_\sigma \bigotimes_{\mathbf{k}\in\mathrm{MBZ}} \gamma_{\mathbf{k}\sigma}^{(-)\dagger} \ket{\Omega_{\mathbf{k}\sigma}}
\end{equation}
The AF phase can be thought of as a long-range coordination between fermionic states related by a $\bm{\pi}$ shift in reciprocal space, producing a new free gas where all interactions have been ``condensed'' in these AF pairs.

\subsection{Spin degeneracy removal}

Up to now, we wrote explicitly the $\sigma$ index. However,
\[
	\epsilon_{\mathbf{k}\uparrow} = \epsilon_{\mathbf{k}\downarrow}
	\qquad
	\Delta_{\mathbf{k}\uparrow} = -\Delta_{\mathbf{k}\downarrow}
\]
thus $E_{\mathbf{k}\uparrow} = E_{\mathbf{k}\downarrow}$. Since the gap function sign does has not a measurable impact, we might as well take the spin index as a pure degeneracy and from now on work in the $\sigma=\uparrow$ sector, dropping the $\sigma$ index. We can further simplify by dropping all indices from the gap function,
\[
	\Delta_{\mathbf{k}\uparrow} = mU \equiv \Delta
\]
being the latter purely $s$-wave. This is a property of the pure HM: for the EHM, the gap acquires a non-trivial topological structure.

\subsection{Chemical potential and system density}

Using the fact that
\[
	\left\langle
		\hat\Gamma_{\mathbf{k}\sigma}^\dagger
		\hat\Gamma_{\mathbf{k}\sigma}
	\right\rangle = \left\langle
		\hat\Psi_{\mathbf{k}\sigma}^\dagger
		\hat\Psi_{\mathbf{k}\sigma}
	\right\rangle
\]
due to the $W_{\mathbf{k}\sigma}$ unitarity properties, the total number of electrons is given by
\begin{align}
	N &= \sum_{\mathbf{k}\sigma} \ev{\hat n_{\mathbf{k}\sigma}} \nonumber \\
	&= \sum_\sigma \sum_{\mathbf{k} \in \mathrm{MBZ}} \left\langle
		\hat n_{\mathbf{k}\sigma} + \hat n_{\mathbf{k}+\bm{\pi}\sigma}
	\right\rangle \nonumber \\
	&= \sum_\sigma \sum_{\mathbf{k} \in \mathrm{MBZ}} \left\langle
		\hat\Gamma_{\mathbf{k}\sigma}^\dagger
		\hat\Gamma_{\mathbf{k}\sigma}
	\right\rangle \nonumber \\
	&= 2 \sum_{\mathbf{k} \in \mathrm{MBZ}} \mathrm{Th}_\mathbf{k}^{(+)}[\beta,\mu] \label{appeq:spinor-particle-number-formula}
\end{align}
being $f$ be the Fermi-Dirac distribution at inverse temperature $\beta$ and chemical potential $\mu$,
\[
	f(\epsilon;\beta,\mu) = \frac{1}{e^{\beta(\epsilon-\mu)}+1} 
\]
The $2$ prefactor appeared because of spin degeneracy. Since the gapped bands refer to the smaller MBZ, thus exhibit the periodicity
\[
	E_\mathbf{k} = E_{\mathbf{k}+\bm{\pi}}
\]
it holds equivalently
\begin{equation}\label{appeq:spinor-particle-number-formula-BZ}
	N = \sum_{\mathbf{k} \in \mathrm{BZ}} \mathrm{Th}_\mathbf{k}^{(+)}[\beta,\mu]
\end{equation}
as is obvious. The $2$ prefactor was absorbed in doubling the integration region, $\mathrm{MBZ}\to\mathrm{BZ}$. Next sections are devoted to simplify the above theoretical results in order to obtain an algorithmic estimation for the magnetization $m$, which is just the AF instability order parameter. 

\begin{figure}
	\centering
	\includegraphics[width=0.8\linewidth]{../Project/HM-HartreeFock/src/labs/n\_AF\_U=10.0\_Δ=1.0.pdf}
	\caption{Electronic density in the AF phase as a function of the chemical potential, measured in relation to the gap $\Delta=mU$, at various temperatures.}
	\label{appfig:af-density-various-betas}
\end{figure}

In Fig.~\ref{appfig:af-density-various-betas} the density $n(\mu)$ is plotted as a function of $\mu$. As is evident, at nearly zero temperature ($\beta\to\infty$) the system is half filled for $-\Delta<\mu<\Delta$. In order to obtain a finite doping, it must be $\abs{\mu}>\Delta$; which means, the chemical potential needs to cross the electronic bands. This HF solution leads invariably to a rather unusual metallic AF at any finite doping.

\subsection{Self-consistency gap equation and magnetization}\label{appsubsec:af-self-consistency-gap-equation}

A convergence algorithm can be designed to find the Hartree-Fock solution to the model. Ultimately, we aim to extract $m$ self-consistently. Combining Eqns.~\eqref{appeq:af-hm-staggered-n-transform}, \eqref{appeq:mixed-cc-term-intermediate-passage} and the definition of Eq.~\eqref{appeq:af-hm-psi-def}, we get
\[
	\frac{1}{2L_xL_y} \sum_\mathbf{r} (-1)^{x+y} \hat n_{\mathbf{r}\sigma} = \sum_{\mathbf{k}\in\mathrm{MBZ}} \left[
		\hat \psi_{\mathbf{k}\sigma} + \hat \psi_{\mathbf{k}\sigma}^\dagger
	\right]
\]
thus
\[
\begin{aligned}
	m &= \frac{1}{2L_xL_y} \sum_\mathbf{r} (-1)^{x+y} \langle
		\hat n_{\mathbf{r}\uparrow} - \hat n_{\mathbf{r}\downarrow}
	\rangle \\
	&= \frac{1}{2L_xL_y} \sum_{\mathbf{k} \in \mathrm{MBZ}} \left\langle
		\left(\hat \psi_{\mathbf{k}\uparrow} + \hat \psi_{\mathbf{k}\uparrow}^\dagger\right) - \left(\hat \psi_{\mathbf{k}\downarrow} + \hat \psi_{\mathbf{k}\downarrow}^\dagger\right)
	\right\rangle
\end{aligned}
\]
Using now Eq.~\eqref{appeq:af-hm-x-psi} we get the self-consistency equation for the HFP $m$,
\begin{align}
	m &= \frac{1}{2L_xL_y} \sum_{\mathbf{k} \in \mathrm{MBZ}} \left[
		\langle
			\hat s_{\mathbf{k}\uparrow}^x
		\rangle
		 - 
		 \langle
			 \hat s_{\mathbf{k}\downarrow}^x
		 \rangle
	\right] \nonumber \\
	&= -\frac{1}{2L_xL_y} \sum_{\mathbf{k} \in \mathrm{MBZ}} \frac{\Re\lbrace\Delta_{\mathbf{k}\uparrow}\rbrace-\Re\lbrace\Delta_{\mathbf{k}\downarrow}\rbrace}{2E_\mathbf{k}} \mathrm{Th}_\mathbf{k}^{(-)}[\beta,\mu] \nonumber \\
	&= -\frac{1}{2L_x L_y} \sum_{\mathbf{k} \in \mathrm{MBZ}} \frac{\Delta}{2E_\mathbf{k}} \mathrm{Th}_\mathbf{k}^{(-)}[\beta,\mu] \label{appeq:af-hm-m-self-consistency}
\end{align}

The magnetization $m$ is the physical order parameter of the phase transition from the normal phase to the antiferromagnetic one. Within such phase transition, in fact, it is the expectation value
\[
	\langle
		\hat c_{\mathbf{k}\sigma}^\dagger \hat 	c_{\mathbf{k}+\bm{\pi}\sigma} 
	\rangle
\]
that goes from zero (normal phase, free Fermi gas) to a non-zero value (AF). As is evident from equations above, $m$ is just its (only) $s$-wave component.

%Now discussion is divided in two parts: zero temperature and finite temperature.
%
%\paragraph{Zero temperature solution.}
%
%For a spin system at zero temperature, the spin operator expectation value anti-aligns with the external field,
%\[
%	\ev{\hat s_{\mathbf{k}\sigma}^x} = \sin(2\theta_{\mathbf{k}\sigma})
%\]
%(see Fig.~\ref{appfig:pseudo-magnetic-field}). Now, since $\Delta_\downarrow = - \Delta_\uparrow$, it follows
%\[
%	\theta_{\mathbf{k}\uparrow} = -\theta_{\mathbf{k}\downarrow} \equiv \theta_\mathbf{k}
%\]
%Then, from Eq.~\eqref{appeq:antiferromagnet-magnetization-self-consistence-abstract}
%\[
%\begin{aligned}
%	m &= \frac{1}{L_xL_y} \sum_{\mathbf{k} \in \mathrm{MBZ}} \sin(2\theta_\mathbf{k}) \\
%	&= \frac{1}{2L_xL_y} \sum_{\mathbf{k} \in \mathrm{BZ}} \sin(2\theta_\mathbf{k})
%\end{aligned}
%\]
%The last passage is due to the fact that the sum over the $\mathrm{MBZ}$ can be performed identically over $\mathrm{BZ}\setminus\mathrm{MBZ}$ and yield the same result. This is because of the lattice periodicity in reciprocal space. Then, finally:
%\begin{equation}\label{appeq:antiferromagnet-magnetization-self-consistence-zero-temperature}
%	m = \frac{1}{L_xL_y} \sum_{\mathbf{k} \in \mathrm{BZ}} [ W_{\mathbf{k}\uparrow} ]_{11} [ W_{\mathbf{k}\uparrow}^\dagger ]_{21}
%\end{equation}
%where $\sin(2\theta_\mathbf{k}) = 2 \sin\theta_\mathbf{k} \cos\theta_\mathbf{k}$ has been used. As will become clear in next section, the sudden appearance of matrix elements of $W$ is not casual.
%
%\paragraph{Finite temperature solution.}
%
%At finite temperature $\beta$, discussion is analogous to the section above. Here will be treated somewhat more theoretically. Define the generic order parameter:
%\[
%	\Delta_{ij} (\mathbf{k}\sigma) \equiv \left\langle
%		[
%			\hat \Psi_{\mathbf{k}\sigma}^\dagger
%		]_i [
%			\hat \Psi_{\mathbf{k}\sigma}
%		]_j
%	\right\rangle
%\]
%In last section, the relevant indices $(i,j)$ were $(1,2)$ and $(2,1)$. Transform this order parameter,
%\begin{align}
%	\Delta_{ij} (\mathbf{k}\sigma) &= \sum_{i'j'} \left\langle
%		[
%			\hat\Gamma_{\mathbf{k}\sigma}^\dagger
%		]_{i'} 
%		[
%			W_{\mathbf{k}\sigma}
%		]_{i'i} [
%			W_{\mathbf{k}\sigma}^\dagger
%		]_{jj'}
%		[
%			\hat\Gamma_{\mathbf{k}\sigma}
%		]_{j'}
%	\right\rangle \nonumber \\
%	&= 
%	\sum_{i'j'} [
%		W_{\mathbf{k}\sigma}
%	]_{i'i} [
%		W_{\mathbf{k}\sigma}^\dagger
%	]_{jj'}
%	\left\langle
%		[
%			\hat\Gamma_{\mathbf{k}\sigma}^\dagger
%		]_{i'} 
%		[
%			\hat\Gamma_{\mathbf{k}\sigma}
%		]_{j'}
%	\right\rangle \nonumber \\
%	&= 
%	\sum_{i'j'} [
%		W_{\mathbf{k}\sigma}
%	]_{i'i} [
%		W_{\mathbf{k}\sigma}^\dagger
%	]_{jj'}
%	\delta_{i'j'} f\left(
%		[d_{\mathbf{k}\sigma}]_{i'i'}; \beta,\mu
%	\right) \nonumber \\
%	&= 
%	\sum_{\ell=1}^2 [
%		W_{\mathbf{k}\sigma}
%	]_{\ell i} [
%		W_{\mathbf{k}\sigma}^\dagger
%	]_{j\ell}
%	f\left(
%		(-1)^\ell E_{\mathbf{k}\sigma}; \beta,\mu
%	\right) \nonumber \\
%	&= 
%	[
%		W_{\mathbf{k}\sigma}
%	]_{1 i} [
%		W_{\mathbf{k}\sigma}^\dagger
%	]_{j 1} f\left(
%		-E_\mathbf{k}; \beta,\mu
%	\right) + [
%		W_{\mathbf{k}\sigma}
%	]_{2 i} [
%		W_{\mathbf{k}\sigma}^\dagger
%	]_{j 2} f\left(
%		E_\mathbf{k}; \beta,\mu
%	\right) \label{appeq:finite-temperature-order-parameter-derivation-intermediate}
%\end{align}
%In the second passage this distribution appeared because an expectation value over a gas of free $\Phi$ fermions was taken. Such an expectation value admits no off-diagonal values, hence the $\delta_{i'j'}$. The diagonal entries are precisely the definition of the Fermi-Dirac distribution for the given energy. At zero temperature and half-filling, the lower band $-E_\mathbf{k}$ is completely filled while the upper band $E_\mathbf{k}$ is empty. Substituting $f=1$ in the first term of line \eqref{appeq:finite-temperature-order-parameter-derivation-intermediate}, $f=0$ in the second and summing $\Delta_{12}$ and $\Delta_{21}$ as done in Eq.~\eqref{appeq:antiferromagnet-magnetization-self-consistence-abstract}, it's easy to derive the result of Eq.~\eqref{appeq:antiferromagnet-magnetization-self-consistence-zero-temperature}. At finite temperature, following the lead of the above paragraph, the antiferromagnetic instability order parameter $m$ will be given simply by
%\begin{equation}\label{appeq:antiferromagnet-magnetization-self-consistence-finite-temperature}
%	m = \frac{1}{2L_xL_y} \sum_{\mathbf{k} \in \mathrm{BZ}} \sum_{\ell=1}^2 
%	[
%		W_{\mathbf{k}\uparrow}
%	]_{\ell 1} [
%		W_{\mathbf{k}\uparrow}^\dagger
%	]_{2 \ell}
%	f\left(
%		(-1)^\ell E_\mathbf{k}; \beta,\mu
%	\right)
%\end{equation}
%Then: mean-field approximations yield an estimate for the magnetization at a given temperature and chemical potential just by carefully combining the elements of the diagonalizing matrix $W$ of each Bogoliubov matrix $h_{\mathbf{k}\sigma}$.

Consider specifically the half-filled situation and multiply Eq.~\eqref{appeq:af-hm-Delta-self-consistency} on both sides by $U$. Since due to symmetry $\mu=0$, we have the finite-temperature self-consistent equation
\[
	\Delta = \frac{U}{2L_xL_y} \sum_{\mathbf{k} \in \mathrm{MBZ}} \frac{\Delta}{\sqrt{\epsilon_{\mathbf{k}}^2 + \Delta^2}} \tanh(\frac{\beta}{2} \sqrt{\epsilon_{\mathbf{k}}^2 + \Delta^2})
\]
where we used the relation
\[
	\mathrm{Th}_\mathbf{k}^{(-)}[\beta,0] = -\tanh(\frac{\beta E_\mathbf{k}}{2})
\]
In thermodynamic limit $L \to \infty$
\[
	\frac{1}{U} \simeq \frac 1 2 \int_\mathrm{MBZ} \frac{\mathrm{d}\mathbf{k}}{(2\pi)^2} \frac{1}{\sqrt{\epsilon_{\mathbf{k}}^2 + \Delta^2}} \tanh(\frac{\beta}{2} \sqrt{\epsilon_{\mathbf{k}}^2 + \Delta^2})
\]
Substitute the momentum integral with an integral over energies, in the low temperature limit, $\beta\to\infty$ (thus approximate the hyperbolic tangent with $1$)
\[
	\frac{1}{U} \simeq \frac 1 2 \int_\mathrm{bands}  \frac{\mathrm{d}\epsilon \rho(\epsilon)}{\sqrt{\epsilon^2 + \Delta^2}}
\]
\todo

\subsection{Instability of the commensurate AF solution and atypical metallic bands}\label{appsubsec:instability-af-commensurate}

A key point to be discussed is the known instability of the AF HF solution of the EHM. As is thoroughly discussed in \cite{singh1990collective}, the solution we presented above is only stable at half-filling, becoming hardly unstable at infinitesimal doping. This is due to the commensurate nature of the established AF phase: we impose SSB to take place at wavevector $\bm{\pi}$, shrinking the BZ to exactly half its size and giving back on the newly defined MBZ two bands $\pm E_\mathbf{k}$. This procedure actually works well for the half-filled case: the magnetization essentially depends on the scattering rate between points on the FS connected by a $\bm{\pi}$ vector. The half filled case, taking advantage of the natural nesting property of the bare bands,
\[
 	\forall \mathbf{k}\in\partial\mathrm{MBZ}\,\colon\, \epsilon_{\mathbf{k}} = \epsilon_{\mathbf{k}+\bm{\pi}}
\]
(being $\partial\mathrm{MBZ}$ the boundary of the MBZ) exhibits a magnetization at infinitesimal interaction. For a many body state resulting from continuous SSB, to establish a given phase means to show the presence of a stable gapless Goldstone mode associated to the broken symmetry. Such a Goldstone mode is essentially a sharp peak at zero frequency for some dressed response function. For the AF phase, it follows from the reduced $\mathrm{U}^z(1)$ rotational symmetry and is related to the proper transverse pseudospin-pseudospin response function,
$2\times 2$ matrix equation,
\begin{equation}\label{appeq:af-commensurate-proper-response}
	\tilde{\chi}_{\hat{\mathbf{S}}^- \hat{\mathbf{S}}^+}(\mathbf{k},\omega) = \frac{\chi_{\hat{\mathbf{S}}^- \hat{\mathbf{S}}^+}(\mathbf{k},\omega)}{1 - U \chi_{\hat{\mathbf{S}}^- \hat{\mathbf{S}}^+}(\mathbf{k},\omega)}
\end{equation}
being $\chi$ the free transverse response. The pseudospin hereby considered is the one of Eq.~\eqref{af-hm-spin-h}. As is shown in \cite{singh1990collective}, for the commensurate AF SSB at wavevector $\bm{\pi}$ the condition for the presence of a stable gapless mode reduces to a $2\times 2$ matrix equation,
\begin{equation}\label{appeq:af-commensurate-stability-condition}
	\text{Stability condition:}
	\qquad
	\mathbb{I}_{2\times2} - U \left[\chi_{\hat{\mathbf{S}}^- \hat{\mathbf{S}}^+}(\bm{\pi},0)\right] \ge 0
\end{equation}
with $\left[\chi_{\hat{\mathbf{S}}^- \hat{\mathbf{S}}^+}(\bm{\pi},\omega)\right]$ the response function matrix obtained by considering a matrix response function for a two-sites unit cell. The above equation is to be interpreted as a matricial equation {\color{tabred}[Explain better..]}. Considering the eigenvalue of largest amplitude of the matrix $\lambda_\delta$ at filling $\delta$, the one most prone to instabilities, it is found that
\[
\begin{aligned}
	\lambda_0(\mathbf{k}\to\bm{\pi},0) &= \frac 1 U &&\text{(half filling)} \\
	\lambda_\delta(\mathbf{k}\to\bm{\pi},0) &= \frac 1 U + \frac{2\epsilon_\mathrm{F}}{\Delta} N(\epsilon_\mathrm{F}) &&\text{(finite doping)}
\end{aligned}
\]
The reason for the distinction above is the intrinsic metallic nature of the system for infinitesimal doping. At half-filling the system is an insulator, therefore no Fermi energy is defined, as well as no FS. Evidently this implies
\[
\begin{aligned}
	1 - U\lambda_0(\mathbf{k}\to\bm{\pi},0) &= 0 &&\text{(half filling)} &&\implies &&\text{stable Goldstone mode}\\
	1 - U\lambda_\delta(\mathbf{k}\to\bm{\pi},0) &< 0 &&\text{(finite doping)} &&\implies &&\text{unstable Goldstone mode}
\end{aligned}
\]
For any finite doping, the commensurate AF phase becomes unstable and is destroyed by transverse spin fluctuations. Nevertheless, for the sake of completeness and because we noticed too late, we will study the metallic AF also at finite doping.

\subsection{Hartree-Fock algorithm results in reciprocal space}\label{appsubsec:hartree-fock-algorithm}

The above sections offer a self-consistency equation for the only HFP $m$; $W$ is determined by $\Delta = mU$ indeed. In an iterative HF fashion, Eq.~\eqref{appeq:af-hm-m-self-consistency} reads:
\begin{equation}\label{appeq:af-hm-m-self-consistency-iterative}
	m_{i+1} =
	\frac{U}{2L_x L_y} \sum_{\mathbf{k} \in \mathrm{MBZ}} \frac{m_i}{2E_\mathbf{k}[\Delta_i]} \left[
	f(-E_\mathbf{k}[m_i];\beta,\mu)-f(E_\mathbf{k}[m_i];\beta,\mu)
	\right]
\end{equation}
Then a self-consistent algorithm to search for a self consistent estimate of the magnetization may be sketched, following the general idea of Sec.~\ref{sec:computational-mft}:
\begin{enumerate}
	\item Algorithm setup: initialize a counter $i=1$ and choose:
	\begin{itemize}
		\item the local repulsion $U/t$ or equivalently the hopping amplitude $t/U$;
		\item the coarse-graining of the $\mathrm{BZ}$, which is, fix $L_x$ and $L_y$. Then 
		\[ 
			k_x = 2\pi n_x/L_x
			\qquad
			k_y = 2\pi n_y/L_y 
		\]
		where $-L_j/2 \le n_j \le L_j/2$, $n_j \in \mathbb{Z}$;
		\item the density $n$ or equivalently the doping $\delta \equiv n-0.5$;
		\item the inverse temperature $\beta$;
		\item the maximum number of iterations $p \in \mathbb{N}$; 
		\item the mixing parameter $g \in [0,1]$;
		\item the tolerance parameters $\delta_m, \delta_n \in \mathbb{R}$ (respectively for the HFP $m$ and density);
	\end{itemize}
	\item Select a random starting value $m_i > 0$ for $i=1$;
	\item Find the optimal chemical potential $\mu$ as follows:
	\begin{enumerate}
		\item Compute $\pm E_\mathbf{k}[\Delta_i, U$]; 
		\item Define the function of the chemical potential $\mu$,
		\[
			\delta n(\mu; \beta, m_i, U) \equiv \frac{1}{L_x L_y} \sum_{\mathbf{k} \in \mathrm{MBZ}} \left[
			f\left(-E_\mathbf{k}[U,m_i]; \beta,\mu \right) + f \left( E_\mathbf{k}[U,m_i]; \beta,\mu \right)
			\right] - n
		\]
		\item Find the root of $\delta n$, 
		\[
			\delta n(\mu_i) = 0
		\]
		Then $\mu_i$ is the chemical potential which, at step $i$ with HFP $m_i$, realizes the best approximation of density $n$;
	\end{enumerate}
	\item Compute $m_{i+1}$ using Eq.~\eqref{appeq:af-hm-m-self-consistency-iterative} with chemical potential $\mu_i$ and update the counter, $i \to i+1$;
	\item Check if $\abs{m_{i+1}-m_i} \le \delta_m$:
	\begin{itemize}
		\item If yes, halt;
		\item If not: check if $i > p$
		\begin{itemize}
			\item If yes, halt and consider choosing better tolerance and model parameters;
			\item If not, define 
			\[
				m_{i+1} \colon g m_{i+1} + (1-g) m_i
			\]
			(logical assignment notation used) and repeat from step 2.
		\end{itemize}
	\end{itemize}
\end{enumerate}

The algorithm sketched hereby was ran with the following setup:
\begin{lstlisting}[language=julia]
UU = [U for U in 0.5:0.5:20.0]      # Local repulsions
LL = [256]             				# Lattice sizes
dd = [d for d in 0.0:0.05:0.45]     # Dopings
bb = [100.0, Inf]				    # Inverse temperatures
p = 100                             # Max number of iterations
dm = 1e-4                           # Tolerance on magnetization
dn = 1e-2                           # Tolerance on density
g = 0.5                             # Mixing parameter
\end{lstlisting}
In Fig.~\ref{appfig:mdU-beta=100} plots at $\beta=100$ of the magnetization $m$ is reported. As is evident in Fig.~\ref{appsubfig:mU-beta=100}, doping disrupts AF ordering; the horizontal asymptote lowers as $\delta$ gets bigger, a mere consequence of the fact that the free space at each single-particle state in $k$-space shrinks as density increases; also, for $\delta > 0$, a finite magnetization $m$ exists only for $U \ge U_c[\delta]$ (a critical value parametrically dependent on doping). Fig.~\ref{appsubfig:dU-beta=100} shows the dependence of magnetization on doping at various fixed local repulsions $U$: once again, to dope the material disrupt AF ordering. Identical analysis have been performed for different temperatures, leading to analogous results as long as $\beta \gtrsim 10$ and to $m \simeq 0$ for higher temperatures (predictably).

\begin{figure}
	\centering
	\subfloat[Local repulsion $U/t$ depencence.]{
		\includegraphics[width=0.5\textwidth]{../Project/HM-HartreeFock/analysis/Phase=AF/scan/Setup=A[128]_t=1.0/PlotOrderParameter/xVar=U/pVar=δ/m_t=1.0_Lx=128.0_β=100.0.pdf}
		\label{appsubfig:mU-beta=100}
	}
	\subfloat[Doping $\delta = n-0.5$ dependence.]{
		\includegraphics[width=0.5\textwidth]{../Project/HM-HartreeFock/analysis/Phase=AF/scan/Setup=A[128]_t=1.0/PlotOrderParameter/xVar=δ/pVar=U/m_t=1.0_Lx=128.0_β=100.0.pdf}
		\label{appsubfig:dU-beta=100}
	}	
	\caption{Plots of the zero-temperature AF instability order parameter $m$ as a function of both the local repulsion $U/t$ (Fig.~\ref{appsubfig:mU-beta=100}) and the doping $\delta = n-1/2$ (Fig.~\ref{appsubfig:dU-beta=100}).}
	\label{appfig:mdU-beta=100}
\end{figure}

%\section{An alternative (less efficient) real-space approach}
%
%The theoretical derivation of the above paragraphs offers a simple description of the system anti-ferromagnetic instability as the instauration of a ground-state of quasiparticles. Here a self-consistent algorithmic extraction of the expected magnetization is presented, following \cite{scholle2023comprehensive}. Note that this algorithm is by far the least efficient, being performed in real space with dimensional exponential scaling in terms of computational time. It is here presented just for completeness as an alternative derivation. This algorithm can become useful for simulations not preserving space-translational invariance, e.g. introducing some degree of disorder. Consider a square lattice of $L_x \times L_y$ sites: the hamiltonian will be a matrix of dimension $2L_x L_y \times 2L_x L_y$,
%\[
%	[ \hat H ]_{(\mathbf{r}\sigma)(\mathbf{r}'\sigma')} = \mel{\Omega}{ \hat c_{\mathbf{r}\sigma} \hat H \hat c_{\mathbf{r}'\sigma'}^\dagger}{\Omega}
%\]
%For simplicity, in the following $D \equiv 2L_x L_y$. In this context, the following convention is used: the rows/column index entry $\alpha=(\mathbf{r}\sigma)$ is associated to a specific site $\mathbf{r}=(x,y)$ and spin $\sigma$ through the relation
%\[
%	\alpha = 2j_\mathbf{r}-\delta_{\sigma=\uparrow}
%	\qq{where}
%	j_\mathbf{r} = x+(y-1)L_x 
%\]
%Let me break this through. For each site $\mathbf{r}$, two sequential indices are provided ($2j_\mathbf{r}-1$, hosting spin $\uparrow$, and $2j_\mathbf{r}$, hosting spin $\downarrow$). $j_\mathbf{r}$ just orders the site rows-wise. This way, $(x,1)$ is assigned to $j_{(x,1)}=x$, while its NN one site above $(x,2)$ is assigned to an entry shifted by $L_x$, $j_{(x,2)}=x+L_x$. This is just a way of counting the site of a finite square lattice by sweeping along a row and then moving to the row above. Fig.~\ref{appfig:square-lattice-ordering} reports a scheme of the used site ordering.
%
%\begin{figure}
%	\centering
%	\input{\GraphicsFolder/square-lattice-ordering}
%	\caption{Schematics of the site ordering on a square lattice performed by sweeping along rows. Left-bottom side is one corner of the lattice. Red sites are characterized by $x+y$ being odd, blue sites by being even. The number reported near to each site is the $\alpha$ entry in the matrix representation for a finite square lattice.}
%	\label{appfig:square-lattice-ordering}
%\end{figure}
%
%Within this convention, matrix elements $H_{\alpha\beta}$ are defined by:
%\begin{itemize}
%	\item If $\sigma=\sigma'$ and $\mathbf{r}$, $\mathbf{r}'$ are NN, the matrix entry is $-t$. In terms of the used greek indices, $\alpha$ and $\beta$ satisfy said requirement if $\abs{\alpha-\beta}=2$ (horizontal hopping) or $\abs{\alpha-\beta}=2L_x$ (vertical hopping). Along column $\alpha$ of the hamiltonian matrix, the elements $-t$ appear at positions 
%	\[
%		( \alpha \pm 2L_x ) \mod D
%		\qq{and}
%		( \alpha \pm 2 ) \mod D
%	\]
%	\item If $\mathbf{r} = \mathbf{r}'$ and $\sigma=\sigma'$ (along the diagonal), the local interaction with the mean field is given by the matrix element 
%	\[
%		-mU \times (-1)^{x+y} \times (-1)^{\delta_{\sigma=\downarrow}}
%	\]
%	Starting from a given entry $\alpha$, $j_\mathbf{r}$ is retrieved simply by $j_\mathbf{r}=\lfloor \alpha/2 \rfloor$, and then
%	\[
%		x+y = (j_\mathbf{r}+1) - \left\lfloor
%			\frac{j_\mathbf{r}}{L_x}
%		\right\rfloor (L_x-1)
%	\]
%	Then the $j_\mathbf{r}$-th $2 \times 2$ block along the diagonal will be given by
%	\[
%		(-1)^{x+y} \underbrace{\begin{bmatrix}
%			-mU & \\ & mU
%		\end{bmatrix}}_{\mathcal{B}}
%	\]
%	Note that the resulting block diagonal contribution to the hamiltonian is shaped like follows (assume $L_x$ to be even):
%	\[
%		\begin{bNiceMatrix}[
%				first-row, 
%				first-col,
%				code-for-first-row = \color{tabred},
%				code-for-first-col = \color{tabred}
%			]
%			 & 1 & 2 & \cdots & L_x-1 & L_x & L_x+1 & L_x+2  & \cdots \\
%			 1 & \mathcal{B} \hfil \\
%			 2 & & -\mathcal{B} \hfil \\
%			 \vdots & & & \ddots \hfil \\
%			 L_x-1 & & & & \mathcal{B} \hfil \\
%			 L_x & & & & & - \mathcal{B} \hfil \\
%			 L_x+1 & & & & & & - \mathcal{B} \hfil \\
%			 L_x+2 & & & & & & & \mathcal{B} \hfil \\
%			 \vdots & & & & & & & & \ddots
%		\end{bNiceMatrix}
%	\]
%	Along the same row, on the diagonal the $2\times2$ blocks $\mathcal{B}$ alternate signs; changing row (in the example above, at positions $L_x,L_x+1$), due to the anti-ferromagnetic configuration of local mean-fields, an additional $-1$ is included. If $L_x$ is taken to be odd, the diagonal blocks just alternate signs all the way.
%\end{itemize}
%These prescriptions allow to build from scratch the hamiltonian matrix. After that, diagonalization provides $D$ orthonormal eigenvectors $\mathbf{v}^\ell \in \mathbb{C}^{D \times D}$ with $\ell=1,\cdots,D$, each associated to a precise eigenvalue $\epsilon^\ell \in \mathbb{R}$. At equilibrium, electrons will fill up the energy eigenstates according to,
%\[
%	\langle
%		\hat n_{\mathbf{r}\sigma}
%	\rangle = \sum_{\ell=1}^D \abs{v_\alpha^\ell}^2 f(\epsilon^\ell;\beta,\mu)
%	\qq{where}
%	\alpha = (\mathbf{r}\sigma)
%\]
%For a fixed filling $n = N/D$, the chemical potential must satisfy
%\[
%\begin{aligned}
%	n &= \frac{1}{D} \sum_{\mathbf{r}\sigma} \langle
%		\hat n_{\mathbf{r}\sigma}
%	\rangle \\
%	&= \frac{1}{D} \sum_{\mathbf{r}\sigma} \sum_{\ell=1}^D \abs{v_\alpha^\ell}^2 f(\epsilon^\ell;\beta,\mu) \\
%	&= \frac{1}{D} \sum_{\ell=1}^D f(\epsilon^\ell;\beta,\mu)
%\end{aligned} 
%\]
%since the $\mathbf{v}^\ell$ eigenvectors are orthonormal. The chemical potential for the half-filled model is already known to be
%\[
%	\mu \big|_{n=1/2} = -\frac{U}{2}
%\]
%as evident from Eq.~\eqref{appeq:af-hm-mft-h}. Average magnetization is then given by
%\begin{align}
%	m &= \frac{1}{D} \sum_{\mathbf{r}} (-1)^{x+y} \left[
%		\langle
%			\hat n_{\mathbf{r}\uparrow}
%			\rangle - \langle
%			\hat n_{\mathbf{r}\downarrow}
%		\rangle
%	\right] \nonumber \\
%	&= \frac{1}{D} \sum_{\lambda=1}^{D/2} (-1)^{
%		(\lambda+1) - \left\lfloor
%			\lambda / L_x
%		\right\rfloor (L_x-1)
%	} \sum_{\ell=1}^D \left[
%		\abs{v_{2\lambda-1}^\ell}^2 - \abs{v_{2\lambda}^\ell}^2
%	\right] f(\epsilon^\ell;\beta,\mu) \label{appeq:average-magnetization-hf-expression}
%%	&= \frac{1}{D} \sum_{\ell=1}^D \left[
%%		1 - 2 \sum_{\lambda=1}^{D/2} \abs{v_{2\lambda}^\ell}^2
%%	\right] f(\epsilon^\ell;\beta,\mu) \nonumber \\
%%	&= n - \frac{2}{D} \sum_{\ell=1}^D \sum_{\lambda=1}^{D/2} \abs{v_{2\lambda}^\ell}^2 f(\epsilon^\ell;\beta,\mu) 
%\end{align}
%since $\mathbf{r}\uparrow$ is associated to an odd index entry, while $\mathbf{r}\downarrow$ to the following even entry. The associated $\mathrm{HF}$ algorithm is identical to the one presented in Sec.~\ref{appsubsec:hartree-fock-algorithm}, with the following substitutions:
%\begin{enumerate}
%	\item[2.] Initialize the hamiltonian matrix $H_{\alpha\beta}$ matrix according to the initialized $m_0$ and the site indexing rules of Fig.~\ref{appfig:square-lattice-ordering};
%	\item[4.] Diagonalize the matrix associated to the operator $\hat H_{\alpha\beta}-\mu \hat N$ collecting the $\mathbf{v}^\ell$ eigenvectors;
%	\item[5.] Compute $m$ using Eq.~\eqref{appeq:average-magnetization-hf-expression} and update the counter, $i \to i+1$;
%\end{enumerate}

\section{(Impossibility of a stable) superconducting phase}\label{appsec:superconducting-phase}

Let us now move to the discussion of superconductivity in the HM. As for the AF phase, we use the results and analyses of this section as groundwork for the more sophisticated case of SC phase in the EHM. We anticipate here that BCS-like superconductivity cannot develop in the pure HM, due to the inherently repulsive nature of $\hat H_U$. BCS theory admits, for an half filled model, finite Cooper fluctuations for an infinitesimal attractive interaction. However, as will be discussed, by the means of MFT it is perfectly possible and coherent to derive a self-consistent superconducting solution to the HM. The point here is that such self-consistent solution sets to zero Cooper fluctuations.

As done for the AF phase in Sec.~\ref{appsec:antiferromagnetic-phase} and following the introductory discussion of Sec.~\mref, let us start from the discussion of the phase symmetries. We aim to describe a uniform-density superconductor,
\[
	\langle \hat n_{i\sigma} \rangle = n
\]
taking into account Cooper RWCs,
\[
	\langle 
		\hat c_{i\sigma}^\dagger 
		\hat c_{i\overline{\sigma}}^\dagger 
	\rangle
\]
which, having we broken $U^\mathrm{c}(1)$ charge symmetry, we allow to fluctuate to non-zero values. For the pure HM, we know from Tab.~\ref{tab:U-wick-terms-phases} that both Hartree and Cooper pairing channels contribute to the RWCs set. We handle them separately.

\subsection{MFT treatment of the Hartree channel}\label{appsubsec:sc-hartree-channel}

For the Hartree channel, knowing the deal with a uniform-density phase, we can re-use the derivation of Sec.~\ref{appsubsec:af-hartree-channel} simply by setting $m=0$. We immediately get from Eq.~\eqref{appeq:af-hm-mft-h}:
\begin{align}
	\hat H_U &\simeq U \sum_\mathbf{r} \left[
		\hat n_{\mathbf{r}\uparrow} \ev{\hat n_{\mathbf{r}\downarrow}} + \ev{\hat n_{\mathbf{r}\uparrow}} \hat n_{\mathbf{r}\downarrow} - \ev{\hat n_{\mathbf{r}\uparrow}} \ev{\hat n_{\mathbf{r}\downarrow}}
	\right] + \text{(Cooper channel)} \nonumber \\
	&= nU \times \hat N - E_{\mathrm{H}/U}^{(\mathrm{SC})} + \text{(Cooper channel)} \label{appeq:sc-hm-U-int-1}
\end{align}
thus chemical potential renormalization $\tilde{\mu} \equiv \mu - nU$ is present also for the SC phase. The contraction energy comes immediately from Eq.~\eqref{appeq:af-hm-hartree-U-Ec},
\begin{equation}\label{appeq:sc-hm-hartree-U-Ec}
	E_{\mathrm{H}/U}^{(\mathrm{SC})} = L_x L_y \times n^2 U
\end{equation}
We now move to Cooper channel.

\subsection{MFT treatment of the Cooper channel}\label{appsubsec:sc-cooper-channel}

Recall the result of Eq.~\eqref{eq:reciprocal-space-local-interaction-explicit} which gave the Fourier-transformed version of the local repulsion, and perform a Cooper class Wick's contraction on it,
\begin{multline}
	\hat H_U \simeq \text{(Hartree channel)} + \frac{U}{L_x L_y}
	\sum_{\mathbf{K}, \mathbf{k}, \mathbf{k}'} \Big[
		\big\langle
			\hat c_{\mathbf{K}+\mathbf{k} \uparrow}^\dagger \hat c_{\mathbf{K}-\mathbf{k} \downarrow}^\dagger
		\big\rangle \hat c_{\mathbf{K}-\mathbf{k}' \downarrow} \hat c_{\mathbf{K}+\mathbf{k}' \uparrow} \\
		+ \hat c_{\mathbf{K}+\mathbf{k} \uparrow}^\dagger \hat c_{\mathbf{K}-\mathbf{k} \downarrow}^\dagger \big\langle
			\hat c_{\mathbf{K}-\mathbf{k}' \downarrow} \hat c_{\mathbf{K}+\mathbf{k}' \uparrow} 
		\big\rangle - \big\langle
			\hat c_{\mathbf{K}+\mathbf{k} \uparrow}^\dagger \hat c_{\mathbf{K}-\mathbf{k} \downarrow}^\dagger
		\big\rangle \big\langle 
			\hat c_{\mathbf{K}-\mathbf{k}' \downarrow} \hat c_{\mathbf{K}+\mathbf{k}' \uparrow}
		\big\rangle
	\Big] \label{appeq:sc-hm-U-int-2}
\end{multline}
We are not breaking translational invariance, thus only Cooper fluctuations with net zero total momentum are allowed. This means only $\mathbf{K}=\mathbf{0}$ contributes. Define as in Eq.~\eqref{eq:su-ehm-phi-def} the pairing operator 
\begin{equation}\label{appeq:sc-hm-phi-def}
	\hat \phi_{\mathbf{k}\uparrow} \equiv \hat c_{-\mathbf{k}\downarrow} \hat c_{\mathbf{k} \uparrow}
\end{equation}
while $\hat n_{\mathbf{k}\sigma}$ is defined as in Eq.~\eqref{appeq:af-hm-n-def}. Eq.~\eqref{appeq:sc-hm-U-int-1} reduces to
\begin{equation}\label{appeq:sc-hm-U-int-3}
	\hat H_U \simeq \text{(Hartree channel)} + \frac{U}{L_x L_y} \sum_{\mathbf{q},\mathbf{q}'} \left[
		\big\langle
			\hat \phi_\mathbf{q}^\dagger
		\big\rangle \hat \phi_{\mathbf{q}'} + 
		\hat \phi_\mathbf{q}^\dagger \big\langle
			\hat \phi_{\mathbf{q}'}
		\big\rangle
	\right] - E_{\mathrm{C}/U}^{(\mathrm{SC})}
\end{equation}
being $E_{\mathrm{C}/U}^{(\mathrm{SC})}$ the HM Cooper contraction energy,
\begin{align}
	E_{\mathrm{C}/U}^{(\mathrm{SC})} &\equiv \frac{U}{L_x L_y} \sum_\mathbf{r} \big\langle
		\hat \phi_\mathbf{q}^\dagger
	\big\rangle \big\langle
		\hat \phi_{\mathbf{q}'}
	\big\rangle \nonumber \\
	&\equiv \frac{U}{L_x L_y} \sum_\mathbf{r}
	\abs{w_\mathbf{q}}^2 \label{appeq:sc-hm-hartree-C-Ec}
\end{align}
having we defined $w_\mathbf{q}$ as in Eq.~\eqref{eq:sc-wqsigma-def},
\[
	w_\mathbf{q} \equiv \big\langle
		\phi_\mathbf{q}^\dagger
	\big\rangle
\]

Evidently in Eq.~\eqref{appeq:sc-hm-U-int-2}, only the $s$-wave component of $w_\mathbf{q}$ partakes,
\begin{equation}\label{appeq:sc-hm-ws-self-consistency}
	w^{(s)} \equiv \frac{1}{L_x L_y} \sum_{\mathbf{q}\in\mathrm{BZ}} w_\mathbf{q}
\end{equation}
thus, defining $\Delta^{(s)} \equiv -U w^{(s)}$, we finally reduce $\hat H_U$ to
\begin{equation}\label{appeq:sc-hm-U-int-4}
	\hat H_U \simeq \text{(Hartree channel)} - \sum_{\mathbf{q}\in\mathrm{BZ}} \left[
	\Delta^{(s)} \hat \phi_\mathbf{q} + 
	\Delta^{(s)*} \hat \phi_\mathbf{q}^\dagger
	\right] - E_{\mathrm{C}/U}^{(\mathrm{SC})}
\end{equation}
which can be easily mapped on a pseudospin model.

\subsection{Nambu representation of the SC phase and diagonalization}\label{appsubsec:sc-nambu}

Composing Eqns.~\eqref{appeq:sc-hm-U-int-1} and \eqref{appeq:sc-hm-U-int-4}, we get the full MFT hamiltonian
\begin{multline}
	\hat H \simeq \sum_{\mathbf{k} \in \mathrm{BZ}} \sum_\sigma \epsilon_{\mathbf{k}\sigma}
	\left[
		\hat c_{\mathbf{k}\sigma}^\dagger \hat c_{\mathbf{k}\sigma}
		- \hat c_{\mathbf{k}+\bm{\pi}\sigma}^\dagger \hat c_{\mathbf{k}+\bm{\pi}\sigma}
	\right] \\
	- \sum_{\mathbf{q}\in\mathrm{BZ}} \left[
	\Delta^{(s)} \hat \phi_\mathbf{q} + 
	\Delta^{(s)*} \hat \phi_\mathbf{q}^\dagger
	\right] + nU \times \hat N - \left[
		E_{\mathrm{H}/U}^{(\mathrm{SC})} + E_{\mathrm{C}/U}^{(\mathrm{SC})}
	\right]
	\label{appeq:sc-hm-mft-h} 
\end{multline}

We now proceed precisely as in Sec.~\ref{appsubsec:af-nambu}. Define the SC Nambu spinor
\[
	\hat \Phi_\mathbf{k} \equiv \begin{bmatrix}
		\hat c_{\mathbf{k}\uparrow} \\
		\hat c_{-\mathbf{k}\downarrow}^\dagger
	\end{bmatrix}
\] 
Define the pseudospin operator:
\[
	\hat s_\mathbf{k}^\alpha \equiv \hat \Phi_\mathbf{k}^\dagger \tau^\alpha \hat \Phi_\mathbf{k}
\]
and the shifted bare bands,
\[
	\xi_\mathbf{k} \equiv \epsilon_\mathbf{k} - \tilde{\mu}
\]
Notice we are using the chemical potential RMP. We employ from the start the spin-independency of the bare bands, as well as their pure $s^*$-wave structure, since no Fock interaction in the present context chan change that. The following chain holds:
\begin{align}
	\hat n_{\mathbf{k}\uparrow} + \hat n_{-\mathbf{k}\downarrow} &= \hat c_{\mathbf{k}\uparrow}^\dagger \hat c_{\mathbf{k}\uparrow} + \hat c_{-\mathbf{k}\downarrow}^\dagger \hat c_{-\mathbf{k}\downarrow} \nonumber \\
	&= \hat c_{\mathbf{k}\uparrow}^\dagger \hat c_{\mathbf{k}\uparrow} - \hat c_{-\mathbf{k}\downarrow} \hat c_{-\mathbf{k}\downarrow}^\dagger + \hat \Id_\mathbf{k} \nonumber \\
	&=  \hat \Phi_\mathbf{k}^\dagger \tau^z \hat \Phi_\mathbf{k} + \hat \Id_\mathbf{k} \nonumber \\
	&=  \hat s_\mathbf{k}^z + \hat \Id_\mathbf{k} \label{appeq:sc-hm-n-z}
\end{align}
Then, considering the kinetic part of the gran-canonical hamiltonian,
\[
\begin{aligned}
	\hat H_t - \tilde{\mu}\hat N &= \sum_\sigma \sum_{\mathbf{k} \in \mathrm{BZ}} \left(
		\epsilon_{\mathbf{k}\sigma} - \tilde{\mu}
	\right) \hat n_{\mathbf{k}\sigma} \\
	&= \sum_{\mathbf{k} \in \mathrm{BZ}} \left[ 
		\xi_\mathbf{k} \hat n_{\mathbf{k}\uparrow} + \xi_{-\mathbf{k}} \hat n_{-\mathbf{k}\downarrow} 
	\right] \\
	&= \sum_{\mathbf{k} \in \mathrm{BZ}} \xi_\mathbf{k} \hat s_\mathbf{k}^z - E_{\mathrm{A}/t}^{(\mathrm{SC})}
\end{aligned}
\]
being $E_{\mathrm{A}/t}^{(\mathrm{SC})}$ the ``anticommutation energy'' coming from the kinetic term,
\begin{equation}
	E_{\mathrm{N}/t}^{(\mathrm{SC})} \equiv -\sum_{\mathbf{k}\in\mathrm{BZ}} \xi_\mathbf{k}
\end{equation}
The scalar term in the above equation is essential for estimating correctly the phase ground state free energy density. The full MFT gran-canonical hamiltonian in Nambu representation is obtained including Eq.~\eqref{appeq:sc-hm-U-int-3},
\begin{equation}\label{sc-hm-bogoliubov-h}
	\hat H - \mu \hat N \simeq \sum_{\mathbf{k} \in \mathrm{BZ}} \sum_\sigma \hat \Phi_\mathbf{k}^\dagger h_\mathbf{k} \hat \Phi_\mathbf{k} - \left[
		E_{\mathrm{A}/t}^{(\mathrm{SC})} + E_{\mathrm{H}/U}^{(\mathrm{SC})} + E_{\mathrm{C}/U}^{(\mathrm{SC})}
	\right]
\end{equation}
being
\[
	h_\mathbf{k} \equiv \begin{bmatrix}
		\xi_\mathbf{k} & -\Delta_\mathbf{k} \\
		-\Delta_\mathbf{k}^* & - \xi_\mathbf{k}
	\end{bmatrix}
\]
while the off-diagonal gap is a constant,
\[
	\Delta_\mathbf{k} \equiv \Delta^{(s)}
\]
This hamiltonian behaves as an ensemble of spins in local magnetic pseudofields
\begin{equation} \label{appeq:sc-hm-pseudofield}
	\mathbf{b}_\mathbf{k} \equiv \begin{bmatrix}
		-\Re{\Delta_\mathbf{k}} \\
		-\Im{\Delta_\mathbf{k}} \\ 
		\xi_\mathbf{k}
	\end{bmatrix}
	\qquad
	\tan 2\theta_\mathbf{k} \equiv \frac{\xi_\mathbf{k}}{\sqrt{\xi_\mathbf{k}^2 + \abs{\Delta_\mathbf{k}}^2}}
	\qquad
	\tan 2\zeta_\mathbf{k} \equiv \frac{\Im{\Delta_\mathbf{k}}}{\Re{\Delta_\mathbf{k}}}
\end{equation}
precisely as was for the AF phase, with the essential difference that now the chemical potential RMP lies inside the $z$ component of the pseudofield itself. This is a consequence of the Nambu spinors definition. We have the spin hamiltonian
\begin{equation}\label{af-su-spin-h}
	\hat H -\mu \hat N \simeq \sum_{\mathbf{k} \in \mathrm{BZ}} \mathbf{b}_\mathbf{k} \cdot \hat{\mathbf{s}}_\mathbf{k} - \left[
		E_{\mathrm{A}/t}^{(\mathrm{SC})} + E_{\mathrm{H}/U}^{(\mathrm{SC})} + E_{\mathrm{C}/U}^{(\mathrm{SC})}
	\right]
	\qq{where}
	\hat{\mathbf{s}}_\mathbf{k} = \begin{bmatrix}
		\hat s_\mathbf{k}^x \\
		\hat s_\mathbf{k}^y \\
		\hat s_\mathbf{k}^z
	\end{bmatrix}
\end{equation}

As was for the AF phase, to diagonalize the hamiltonian means essentially to align the spin $z$ axis with the direction of the pseudofield. Then, we diagonalize the hamiltonian via a set of rotations similarly to Eq.~\eqref{appeq:af-hm-diagonalizer}. The main difference is that in general the pseudofield can have a $y$ component, thus the diagonalizing matrix is made of a sequence of two rotations, both counter-clockwise, the first around the $z$ axis by an angle $2\zeta_\mathbf{k}$ to place the pseudofield in the $xz$ plane, the second around the new $y$ axis by an angle $2\theta_\mathbf{k}$ to align the $z$ axis with the pseudofield,
\begin{align}
	W_\mathbf{k} &\equiv e^{-i \theta_\mathbf{k} \tau^y} e^{-i \zeta_\mathbf{k} \tau^z} \nonumber \\
	&= \begin{bmatrix}
		\cos \theta_\mathbf{k} & - \sin \theta_\mathbf{k} \\
		\sin \theta_\mathbf{k} & \cos \theta_\mathbf{k}
	\end{bmatrix} \begin{bmatrix}
		e^{-i\zeta_\mathbf{k}} & \\ 
		& e^{i\zeta_\mathbf{k}}
	\end{bmatrix} \nonumber \\
	&= \begin{bmatrix}
		\cos \theta_\mathbf{k} e^{-i\zeta_\mathbf{k}} & - \sin \theta_\mathbf{k} e^{i\zeta_\mathbf{k}} \\
		\sin \theta_\mathbf{k} e^{-i\zeta_\mathbf{k}} & \cos \theta_\mathbf{k} e^{i\zeta_\mathbf{k}}
	\end{bmatrix} \label{appeq:sc-hm-diagonalizer}
\end{align}
obtaining the diagonal matrix
\[
	d_\mathbf{k} \equiv W_\mathbf{k} h_\mathbf{k} W_\mathbf{k}^\dagger 
	\qq{where}
	d_\mathbf{k} \equiv \begin{bmatrix}
		E_\mathbf{k} & \\ & -E_\mathbf{k}
	\end{bmatrix}
\]
and the bands
\[
	E_\mathbf{k} \equiv \sqrt{\xi_\mathbf{k}^2 + \abs{\Delta_\mathbf{k}}^2}
\]

Finally, the Nambu spinor for Bogoliubov quasiparticles is given similarly to Eq.~\eqref{appeq:af-hm-Gamma-WPsi}
\begin{equation}\label{appeq:sc-hm-Gamma-WPsi}
	\hat\Gamma_\mathbf{k} = W_\mathbf{k} \hat \Phi_\mathbf{k}
\end{equation}
Explicitly,
\begin{align}
	\hat\Gamma_\mathbf{k} &= \begin{bmatrix}
		\cos \theta_\mathbf{k} e^{-i\zeta_\mathbf{k}} & - \sin \theta_\mathbf{k} e^{i\zeta_\mathbf{k}} \\
		\sin \theta_\mathbf{k} e^{-i\zeta_\mathbf{k}} & \cos \theta_\mathbf{k} e^{i\zeta_\mathbf{k}}
	\end{bmatrix} \begin{bmatrix}
		\hat c_{\mathbf{k}\uparrow} \\
		\hat c_{-\mathbf{k}\downarrow}^\dagger 
	\end{bmatrix} \nonumber \\
	&= \begin{bmatrix}
		\cos \theta_\mathbf{k} e^{-i\zeta_\mathbf{k}} \hat c_{\mathbf{k}\uparrow} - \sin \theta_\mathbf{k} e^{i\zeta_\mathbf{k}} \hat c_{-\mathbf{k}\downarrow}^\dagger  \\
		\sin \theta_\mathbf{k} e^{-i\zeta_\mathbf{k}} \hat c_{\mathbf{k}\uparrow} + \cos \theta_\mathbf{k} e^{i\zeta_\mathbf{k}} \hat c_{-\mathbf{k}\downarrow}^\dagger
	\end{bmatrix} \equiv
	\begin{bmatrix}
		\hat \gamma_\mathbf{k}^{(+)} \\
		\hat \gamma_\mathbf{k}^{(-)} 
	\end{bmatrix} \label{appeq:sc-hm-gamma-ops-transform}
\end{align}
These operators act as free fermion fields on the Bogoliubov bands $\pm E_\mathbf{k}$.

\subsection{Explicit zero-temperature SC ground state at half-filling}\label{appsubsec:sc-zero-T-half-filling-gs}

Reasoning precisely as in Sec.~\ref{appsubsec:af-zero-T-half-filling-gs}, we can extract the half-filling zero temperature ground state, commonly referred to as the ``BCS state''. Eqns.~\eqref{appeq:af-hm-gs-1} and \eqref{appeq:af-hm-tilted-down} hold equivalently (ignoring the $\sigma$ index and extending momentum tensor product to the entire BZ),
\[
	\ket{\Psi} = \bigotimes_{\mathbf{k}\in\mathrm{BZ}} \left[
		\cos \theta_\mathbf{k} \ket{\Downarrow_\mathbf{k}} - \sin \theta_\mathbf{k} \ket{\Uparrow_\mathbf{k}}
	\right]
\]
Using now Eq.~\eqref{appeq:sc-hm-n-z}, we have
\[
\begin{aligned}
	\mel{\Downarrow_\mathbf{k}}{\hat s_\mathbf{k}^z + \hat \Id_\mathbf{k}}{\Downarrow_\mathbf{k}} &= 0 \\
	\mel{\Uparrow_\mathbf{k}}{\hat s_\mathbf{k}^z + \hat \Id_\mathbf{k}}{\Uparrow_\mathbf{k}} &= 2
\end{aligned}
\]
so up to a phase
\[
	\ket{\Downarrow_\mathbf{k}} \sim \ket{\Omega_\mathbf{k}}
	\qquad
	\ket{\Uparrow_\mathbf{k}} \sim \hat \phi_\mathbf{k}^\dagger \ket{\Omega_\mathbf{k}}
\]
Hence, we have
\[
	\ket{\Psi} = \bigotimes_{\mathbf{k}\in\mathrm{BZ}} \left[
		\cos \theta_\mathbf{k} + e^{2i\zeta_\mathbf{k}} \sin \theta_\mathbf{k} \hat \phi_\mathbf{k}
	\right] \ket{\Omega_\mathbf{k}}
\]
The argument $2\zeta_\mathbf{k}$ of the phase factor is precisely the phase of the gap function, as can be shown easily.

\subsection{Superconducting fluctuations suppression in the conventional HM}\label{appsubsec:sc-hm-impossible}

Let us report the SC analogous of the AF Eqns.~\eqref{appeq:af-hm-x-psi} and \eqref{appeq:af-hm-y-psi}:
\begin{align}
	\hat s_\mathbf{k}^x &= \hat \phi_\mathbf{k} + \hat \phi_\mathbf{k}^\dagger \label{appeq:sc-hm-x-psi} \\
	\hat s_\mathbf{k}^y &= -i \left(
	\hat \phi_\mathbf{k} - \hat \phi_\mathbf{k}^\dagger 
	\right) \label{appeq:sc-hm-y-psi}
\end{align}
while the analogous of Eq.~\eqref{appeq:af-hm-z-psi} was given in Eq.~\eqref{appeq:sc-hm-n-z}. Let also the thermal factor
\begin{equation}\label{appeq:sc-thermal-factors-def}
	\mathrm{Th}_\mathbf{k}^{(\pm)}[\beta] \equiv 
	f\left(
		E_\mathbf{k};\beta,0
	\right) \pm f\left(
		-E_\mathbf{k};\beta,0
	\right)
\end{equation}
and the EVs:
\begin{align}
	\ev{\hat s_\mathbf{k}^x} &= -\frac{\Re\lbrace\Delta_\mathbf{k}\rbrace}{2E_\mathbf{k}} \mathrm{Th}_\mathbf{k}^{(-)}[\beta] \label{appeq:sc-hm-sx}\\
	\ev{\hat s_\mathbf{k}^y} &= -\frac{\Im\lbrace\Delta_\mathbf{k}\rbrace}{2E_\mathbf{k}}  \mathrm{Th}_\mathbf{k}^{(-)}[\beta] \label{appeq:sc-hm-sy}\\
	\ev{\hat s_\mathbf{k}^z} &=  \frac{\xi_\mathbf{k}}{2E_\mathbf{k}}  \mathrm{Th}_\mathbf{k}^{(-)}[\beta] \label{appeq:sc-hm-sz}
\end{align}

Up to now, the derivation was general. However, it is immediate to show that for the pure HM the SC phase can never establish. Recall Eq.~\eqref{appeq:sc-hm-ws-self-consistency}: explicitly
\[
\begin{aligned}
	w^{(s)} &= \frac{1}{L_x L_y} \sum_{\mathbf{k}\in\mathrm{BZ}} \big\langle
		\hat \phi_\mathbf{k}^\dagger
	\big\rangle \\
	&= \frac{1}{L_x L_y} \sum_{\mathbf{k}\in\mathrm{BZ}} \big\langle
		\hat \Phi_\mathbf{k}^\dagger \tau^+ \hat \Phi_\mathbf{k}
	\big\rangle \\
	&= \frac{1}{L_x L_y} \sum_{\mathbf{k}\in\mathrm{BZ}} \left[
		\ev{\hat s_\mathbf{k}^x} + i \ev{\hat s_\mathbf{k}^y}
	\right] \\
	&= - \frac{1}{L_x L_y} \sum_{\mathbf{k}\in\mathrm{BZ}} \frac{\Delta_\mathbf{k}}{E_\mathbf{k}} \mathrm{Th}_\mathbf{k}^{(-)}[\beta]
\end{aligned}
\]
Now: $\Delta_\mathbf{k} = \Delta^{(s)} = -U w^{(s)}$, thus if we define
\[
	U_\mathrm{c}[\beta] \equiv \left[
		\frac{1}{L_x L_y} \sum_{\mathbf{k}\in\mathrm{BZ}} \frac{\mathrm{Th}_\mathbf{k}^{(-)}[\beta]}{E_\mathbf{k}}
	\right]^{-1}
\]
(the reason for this notation will be clear in Sec.~\mref) we end with the equation:
\[
	w^{(s)} = -\frac{U}{U_\mathrm{c}[\beta]} w^{(s)}
\]
Since evidently $U_\mathrm{c}[\beta] > 0$ by definition, the above equation can be solved in $\mathbb{C}$ only by $w^{(s)}=0$. This implication terminates the argument: superconductivity cannot establish self-consistently at this level in the HM.